% !TEX program = xelatex
\documentclass[a4paper]{article}
\usepackage{amsthm}
\usepackage{amssymb}
\usepackage{bm}
\usepackage{mathtools}
\usepackage[x11names]{xcolor}
\usepackage{xparse}
\usepackage{fontspec}
\usepackage{unicode-math}
\setromanfont{DovesType-Regular.otf}
\setsansfont{Andika}
\setmathfont{Asana Math}[Scale=1]

\usepackage{varwidth}
\usepackage{siunitx}
\usepackage{graphicx}
\usepackage[margin=1cm]{geometry}
\usepackage[most]{tcolorbox}
\usepackage{pgfplots}
\pgfplotsset{compat=newest}
\tcbuselibrary{skins,xparse,poster}
% \usetikzlibrary{fadings}
\usetikzlibrary{calc, plotmarks, shapes, shapes.geometric, positioning, angles, intersections, quotes, through, patterns, turtle, arrows.meta}
\usetikzlibrary{decorations.markings}
% \usepackage{etoolbox}
\usepackage{tkz-euclide}

%%%%%%%%%%%%%%%%%%%%%%%%%%%%%%%%%%%%%%%%%%%%%%%%%%%%%%%%%
\newcommand\markangle[9]{% origin X Y radius radiusmark mark colour opacity
%  % fill red circle offset-from-centre
  \begin{scope}
    \path[clip] (#1) -- (#2) -- (#3);
    \fill[color=#7,fill opacity=#8,draw=black,name path global=pcircle]  % global declaration required otherwise pcircle is not seen by the `named intersections=' lines below.
    (#1) circle (#4);
  \end{scope}
  % middle calculation
  \path[name path=line one] (#1) -- (#2);
  \path[name path=line two] (#1) -- (#3);
  \path[%
  name intersections={of=line one and pcircle, by={inter one}},
  name intersections={of=line two and pcircle, by={inter two}}
  ] (inter one) -- (inter two) coordinate[pos=#9] (place);
  % put mark
  \node at ($(#1)!#5!(place)$) {\scriptsize{#6}};
}

% \newcommand{\condSoln}[2]{\ifcsdef{r@#1}{#2}{}}

% \newcommand\fadingtext[3][]{%
%    \begin{tikzfadingfrompicture}[name=fading letter]
%      \node[text=transparent!0,inner xsep=0pt,outer xsep=0pt,#1] {#3};
%    \end{tikzfadingfrompicture}%
%    \begin{tikzpicture}[baseline=(textnode.base)]
%      \node[inner sep=0pt,outer sep=0pt,#1](textnode){\phantom{#3}};
%      \shade[path fading=fading letter,#2,fit fading=false]
%      (textnode.south west) rectangle (textnode.north east);%
%    \end{tikzpicture}%
% }

\definecolor{JISpurple}{RGB}{89,72,122}
\definecolor{JISivory}{RGB}{241,234,221}
\definecolor{JIStaupe}{RGB}{183,156,154}
\definecolor{PaleGreen}{RGB}{240,255,240} % 'Honeydew'

\AddToHook{shipout/background}{%
    \put (0in,-\paperheight){\includegraphics[width=\paperwidth,height=\paperheight]{images/R10V5-R.png}}%
}

\newtcolorbox{MyOuterBox}{%
  enhanced,
  frame style=JISpurple,
  colback=JISivory,
  colframe=JISpurple,
  title={\includegraphics[width=0.9cm,height=0.9cm]{images/JIS Final Logo FA-02.png}\raisebox{3mm}{\Large{Maths Challenge}\hspace{26em} \Large{\bfseries\sffamily 3}}},
}

\newtcolorbox{MyInnerBox}[2][]{enhanced,%empty,
coltitle=JISpurple,colback=white,
fonttitle=\bfseries\sffamily,
attach boxed title to top left={yshift=-1.5mm},
boxed title style={empty, size=small, top=1mm, bottom=0pt},
varwidth boxed title=0.5\linewidth,
frame code={
  \path (title.east|-frame.north) coordinate (aux);
\path[draw=JISpurple, line width=0.5mm, rounded corners,fill=white]
(frame.west) |- ([xshift=-2.5mm]title.north east) to[out=0, in=180] ([xshift=7.5mm]aux)-|(frame.east)|-(frame.south)-|cycle;
},
title={#2},#1}

\newtcolorbox{MyInnerSplitBox}[2][]{enhanced,%empty,
bicolor,sidebyside,sidebyside align=top seam,
righthand width=5.5cm,colbacklower=white,
coltitle=JISpurple,colback=white,
fonttitle=\bfseries\sffamily,
attach boxed title to top left={yshift=-1.5mm},
boxed title style={empty, size=small, top=1mm, bottom=0pt},
varwidth boxed title=0.5\linewidth,
frame code={
  \path (title.east|-frame.north) coordinate (aux);
\path[draw=JISpurple, line width=0.5mm, rounded corners,fill=white]
(frame.west) |- ([xshift=-2.5mm]title.north east) to[out=0, in=180] ([xshift=7.5mm]aux)-|(frame.east)|-(frame.south)-|cycle;
},
title={#2},#1}


\newtcolorbox{MySolutionBox}{%
  enhanced, frame style=JISpurple,
  colback=PaleGreen, colframe=green,
  title={\Large Solution},
  drop fuzzy shadow,
}

%%%%%%%%%%%%%%%%%%%%%%%%%%%%%%%%%%%%%%%%%%%%%%%%%%
\newtoggle{SOLUTION}
%%% Uncomment the appropriate line below to show solutions %%%
% \toggletrue{SOLUTION}
\togglefalse{SOLUTION}
%%%%%%%%%%%%%%%%%%%%%%%%%%%%%%%%%%%%%%%%%%%%%%%%%


%%%%%%%%%%%%%%%%%%%%%%%%%%%%%%%%%%%%%%%%%%%%%%%%%%
%%%%%%            DOCUMENT BEGINS           %%%%%%
%%%%%%%%%%%%%%%%%%%%%%%%%%%%%%%%%%%%%%%%%%%%%%%%%%
\begin{document}


  \begin{MyOuterBox}
    \iftoggle{SOLUTION}{Here are the full, or partial solutions.
    }{
      Welcome to this week's Maths Challenge!\\
      Everyone is encouraged to take part.\\
      There are 3 to 5 merit points for correct solutions, depending on the difficulty of the problem and how impressed the marker is with your solution.\\
      There are house points for students who submit 7 or more correct solutions in one term.\\
      Solutions must be explained in detail, responses with just the answer may be ignored.\\
      Drop your solution in the box in the staffroom by Tuesday.
    }
       \begin{MyInnerSplitBox}{Year 8 and below}
     Six identical rhombuses, each of area \(\SI{5}{\square\cm}\), form a star. The tips of the star are joined to form a regular hexagon, as shown in blue.\par
     What is the area of the hexagon?
      \iftoggle{SOLUTION}{%conditional output begin
      \begin{MySolutionBox}
        It's a good idea to start off by marking some angles on the diagram.\par
        The hexagon is regular so all the internal angles are the same. The rhombuses (sides are all the same length) are identical, so that means the angles at the center, O, are all equal, so each rhombus has an acute angle of \(\SI{60}{\degree} \left(= \frac{\SI{360}{\degree}}{6}\right)\), see the blue angle \(I_{0}\widehat{O}I_{1}\).\par
        Opposite angles in a rhombus are equal and the angles in a quadrilateral sum to \(\SI{360}{\degree}\) so the obtuse angles in each rhombus are \(\SI{120}{\degree}\), see the green angle \(H_{0}\widehat{I_{0}}O\).\par
        At \(I_{n}\) the angles in each of the rhombuses are each \(\SI{120}{\degree}\) so angle \(H_{0}\widehat{I_{1}}H_{1}\) must also be \(\SI{120}{\degree}\). (See the pink angle for example.)\par
      Each of the outer white triangles, for example, \(H_{1}I_{2}H_{2}\) are isosceles because the sides of the rhombuses are equal. This means that all the angles like \(I_{2}\widehat{H_{2}}H_{1}\) must be \(\SI{30}{\degree}\), see the yellow angles.\par
      So each white triangle is exactly half a rhombus and has an area of \(\SI{2.5}{\square\cm}\). There are six of them having a total area of \(6\times 2.5 = \SI{15}{\square\cm}\). Add on the area of the six rhombuses for a total area of \(\SI{45}{\square\cm}\).
      \end{MySolutionBox}
      }{}%conditional output end
      \tcblower
      \begin{tikzpicture}[scale=1,
        decoration={markings, mark= at position 0.5 with {\draw[-] (0,-2pt) -- (0,2pt);}}]
        \newdimen\R
        \R=2.5cm
        \coordinate (O) at (0:0);
        \foreach \x in {0,...,5}{
          \path[draw,fill=yellow!10!white] (0:0) -- ({\x*60-30}:{\R*(1/sqrt(3))}) coordinate (I\x) -- (\x*60:\R) -- ({\x*60+30}:{\R*(1/sqrt(3))});
          \path[draw,blue] ({\x*60}:\R) coordinate (H\x) -- ({(\x+1)*60}:\R);
          \iftoggle{SOLUTION}{
            \path (0:0) decorate { -- ({\x*60-30}:{\R*(1/sqrt(3))}) coordinate (I\x)} decorate { -- (\x*60:\R)} decorate { -- ({\x*60+30}:{\R*(1/sqrt(3))})};
            \node[label={[font=\scriptsize,label distance=-2mm]{\x*60}:{\({H_{\x}}\)}}] at (H\x){};
            \node[label={[font=\scriptsize,label distance=-1.5mm]{\x*60-30}:{\({I_{\x}}\)}}] at (I\x){};
          }{
            \node[blue] at ({\x*60}:1.25) {\Large 5};
          }
        }
        \iftoggle{SOLUTION}{%
          \node[label={[font=\scriptsize,xshift=-0.3cm, yshift=-0.35cm]O}] at (O) {};
          \markangle{O}{I0}{I1}{3mm}{6mm}{\(\SI{60}{\degree}\)}{SteelBlue1}{0.5}{0.5}
          \markangle{I0}{I3}{H0}{3mm}{5mm}{\(\SI{120}{\degree}\)}{DarkSeaGreen2}{0.5}{0.5}
          \markangle{I2}{H2}{H1}{3mm}{5mm}{\(\SI{120}{\degree}\)}{LightPink1}{0.5}{0.5}
          \markangle{H2}{I2}{H1}{3mm}{6mm}{\(\SI{30}{\degree}\)}{Gold1}{0.5}{0.5}
          \markangle{H1}{I2}{H2}{3mm}{5mm}{\(\SI{30}{\degree}\)}{Gold1}{0.5}{0.5}
        }{}
      \end{tikzpicture}
    \end{MyInnerSplitBox}


    \vspace{0.4cm}
        \begin{MyInnerSplitBox}{Year 9 and above}
      In a regular pentagram (5-pointed star), (i) show that the angle at each point is \(\SI{36}{\degree}\).\par
      Thus the sum of the angles in all five points is \(\SI{180}{\degree}\).\par
      Now, (ii) what is the sum of the angles in all five points of an irregular pentagram, as illustrated in the second diagram?
      \iftoggle{SOLUTION}{%conditional output begin
      \begin{MySolutionBox}
        (i) Consider the quadrilateral shown in bold on the regular pentagram. We know that its internal angles should sum to \(\SI{360}{\degree}\) and we know that three of its interior angles are the same, that is, three times the interior angle of a regular pentagon. Recall the sum of interior angles of an n-sided polygon is \(180\times (n-2)\), so for a pentagon,
        \begin{align*}
          180\times 3 = \SI{540}{\degree} \text{, and one angle is } \frac{540}{5} = \SI{108}{\degree}
          \shortintertext{Then the angle at the point}
          = 360 - (3\times 108) = \SI{36}{\degree}
        \end{align*}
        (ii) There is a nice graphical solution to this. Place an arrow along one side of the irregular pentagram. Rotate the arrow around its base, clockwise, so that it lies over the line indicated by the green arrow (1). Imagine that at the same time the red line shrings to the size of the line it lies over. Now, we rotate the red line again clockwise, this time around the arrow tip (see green arrow 2). The next rotation of the red line is shown with green arrow 3. If we repeat this two more times the red arrow will lie over the same line it started on, but, the arrow will be at the other end. In other words, we have turned the red line through all the angles of the pentagram and it has turned \(\SI{180}{\degree}\). So the sum of the angles in all five points is \(\SI{180}{\degree}\)! Isn't that neat?
      \end{MySolutionBox}
    }{}%conditional output end
    \tcblower
    \begin{tikzpicture}[scale=1]
      \pgfmathsetmacro{\side}{5} %length of line
      \pgfmathsetmacro{\ang}{60} %angle of the first line anticlockwise from x-axis
      \draw[blue] (45:{sqrt(2)}) coordinate (R0) -- ++(\ang:\side) coordinate (R1) -- ++({\ang+216}:\side) coordinate (R2) -- ++({\ang+72}:\side) coordinate (R3) -- ++({\ang+288}:\side) coordinate (R4) -- ++({\ang+144}:\side);
      % \draw[green] (1,1) -- ++(\ang:3)  -- ([turn]-144:3) -- ([turn]-144:3) -- ([turn]-144:3) -- ([turn]-144:3); % alternative using 'turn' notation, easier!
      \iftoggle{SOLUTION}{%
      \tkzInterLL(R0,R1)(R2,R3) \tkzGetPoint{Ra}
      \tkzInterLL(R3,R4)(R0,R1) \tkzGetPoint{Rb}
      \tkzInterLL(R1,R2)(R3,R4) \tkzGetPoint{Rc}
      \draw[blue,ultra thick] (R2) -- (Ra) -- (Rb) -- (Rc) -- cycle;
      \markangle{R2}{R3}{R1}{6mm}{8mm}{\(\SI{36}{\degree}\)}{red}{0.5}{0.5}
      \markangle{Ra}{R2}{R1}{3mm}{5.5mm}{\(\SI{108}{\degree}\)}{SteelBlue1}{0.5}{0.5}
      \markangle{Rb}{R0}{R4}{3mm}{5mm}{\(\SI{108}{\degree}\)}{SteelBlue1}{0.5}{0.5}
      \markangle{Rc}{R3}{R2}{3mm}{5.5mm}{\(\SI{108}{\degree}\)}{SteelBlue1}{0.5}{0.5}
    }{}
    \end{tikzpicture}%
    \par\vspace{0.5cm}
    \begin{tikzpicture}[scale=1]
      \pgfmathsetmacro{\side}{4} %length of line
      \pgfmathsetmacro{\ang}{60} %angle of the first line anticlockwise from x-axis
      \draw[] (45:{sqrt(2)}) coordinate (P1)-- ++(\ang:\side) coordinate (P2) -- ++({\ang+225}:{\side+1.5}) coordinate (P3) -- ++({\ang+62}:{\side+1.75}) coordinate (P4) -- ++({\ang+289}:\side) coordinate (P5) -- ++({\ang+147.7}:{\side+0.86});
      \iftoggle{SOLUTION}{
        \draw[ultra thick,red,-Stealth] (P2) -- (P3);
        \coordinate (c) at ($(P2)!0.25!(P3)$);
        \coordinate (d) at ($(c)!1.1cm!-110:(P3)$);
        \draw[-Stealth,green,thick] (c) node[black,below left] {\(1\)} to [bend left] (d);
        \draw[ultra thick,red!60,-Stealth] (P2) -- (P1);
        \coordinate (e) at ($(P1)!0.25!(P2)$);
        \coordinate (f) at ($(e)!0.5cm!250:(P2)$);
        \draw[Stealth-,green,thick] (f) node[black,above right] {\(2\)} to [bend right] (e);
        \draw[ultra thick,red!30,-Stealth] (P5) -- (P1);
        \draw[ultra thick,red!15,-Stealth] (P5) -- (P4);
        \coordinate (g) at ($(P5)!0.25!(P1)$);
        \coordinate (h) at ($(g)!0.7cm!60:(P5)$);
        \draw[-Stealth,green,thick] (g) node[black,above left] {\(3\)} to [bend left] (h);
      }{}
    \end{tikzpicture}
    \end{MyInnerSplitBox}


    %%%%%%%%%%%% 5-pointed star
    % \begin{tikzpicture}[scale=1]
    %   \pgfmathsetmacro{\ct}{3} % distance center to tip
    %   \pgfmathsetmacro{\cc}{\ct*sin(18)/sin(126)} % distance center to corner (sine rule)
    %   % star
    %   \draw[thick,blue] (0,0)
    %       +(90-0*36:\ct) coordinate(T1)
    %       foreach[evaluate=\x as \nc using int((\x+1)/2),   % number for corner coordinates
    %               evaluate=\x as \nt using int((\x+1)/2+1)] % number for tip coordinates
    %           \x in {1,3,...,9}{
    %           -- +(90-\x*36:\cc) coordinate(C\nc) -- +({90-(\x+1)*36}:\ct) coordinate(T\nt)}
    %       -- cycle;
    %     \markangle{T1}{T3}{T4}{6mm}{8mm}{\(\SI{36}{\degree}\)}{SteelBlue1}{0.5}{0.5}
    % \end{tikzpicture}
    %%%%%%%%%%%%%%%% Pentagram (regular)
     % % Turtles!
     %  \tikz[turtle/distance=4.5cm]
     %  \draw[turtle=home,]
     %    \foreach \i in {1,...,5}{
     %      [turtle={forward,right=144}]
     %    }

% Irregular pentagram with turtle
     % \draw[turtle={home,left=40,forward=3cm,right=150,forward=4cm,right=150,forward=3cm,right=110,forward=3cm,right=160,forward=3.4cm}]

  \end{MyOuterBox}

%%%%%%%%%%%%%%%%%%%%%%%%%%%%%%%%%%%%%%%%%%%%%%%%%%
\end{document}



